\weeknum{8}
\topic[Steady-State Heat Flow]{Geothermics: Steady-State Heat Flow}

\workshopstart

\vspace{-0.75em}
\noindent\textbf{Duration:} 3 hours \hfill \textbf{Setting:} Workshop room

\section*{Preparation}

\begin{description}
  \item[Pre-reading:] Course Notes, Ch. 10
  \item[Lecture videos:]
  \item[Notes:]
\end{description}

\section*{Purpose}

This workshop uses conservation-of-energy reasoning to build intuition for steady-state geotherms: how thermal conductivity contrasts refract isotherms and redistribute gradients, and how crustal heat production controls geotherm curvature and deep temperatures.

\section*{Learning Outcomes}
By the end of this workshop, students should be able to:
\begin{itemize}
  \item Use conservation of energy to reason about steady-state geotherms.
  \item Explain how thermal conductivity contrasts redistribute isotherms.
  \item Interpret geotherm curvature in terms of crustal heat production.
  \item Relate conductivity and heat production variations to geothermal potential.
\end{itemize}

\section*{Session Flow (\timelinelength\ min)}

{\color{red}
\timeline[slot]{Warm-up --- Where does Earth's heat come from?}{10}

Ask students to list (in pairs) all the sources of heat they can think of inside the Earth. Collect responses and briefly categorise them: primordial heat (accretion, differentiation), radiogenic decay, and tidal dissipation. This motivates why steady-state geotherms vary between terranes and why heat production is a key parameter.
}

\timeline[slot]{Mini-lecture --- Conceptual Framing}{15}
\begin{itemize}
  \item Steady-state heat flow as an energy bookkeeping problem.
  \item Review Fourier's law and the meaning of heat flow.
  \item Emphasize the separation of roles:
  \begin{itemize}
    \item conductivity controls \emph{gradients},
    \item heat production controls \emph{curvature}.
  \end{itemize}
\end{itemize}

\timeline[item]{Refraction Exercise I: Insulating Sediments}{30}
\paragraph{Setup}
Compare two crustal columns with identical basal heat flow:
\begin{itemize}
  \item Case A: crystalline crust exposed at the surface.
  \item Case B: crystalline crust overlain by a low-conductivity sedimentary layer.
\end{itemize}

\paragraph{Activities}
\begin{itemize}
  \item Predict temperature gradients in sediments vs basement.
  \item Sketch qualitative isotherms for each case.
  \item Identify where isotherms compress and where they stretch.
\end{itemize}

\paragraph{Key Concepts}
\begin{itemize}
  \item Heat flow conserved across layers.
  \item Lower conductivity $\rightarrow$ steeper gradient.
  \item Sediments act as a thermal blanket.
\end{itemize}

\paragraph{Geological Discussion}
\begin{itemize}
  \item Implications for geothermal gradients measured in boreholes.
  \item Relative geothermal potential of sedimentary basins vs exposed cratons.
\end{itemize}

\timeline[item]{Refraction Exercise II: Layered Crustal Materials}{35}
\paragraph{Focus}
Thermal refraction arising from conductivity contrasts within the crust.

\paragraph{Examples}
\begin{itemize}
  \item Shales (low $k$) vs quartzites (high $k$).
  \item Cyclic sedimentary layering preserved through metamorphism.
\end{itemize}

\paragraph{Activity}
\begin{itemize}
  \item Apply $q = -k \nabla T$ across layered media.
  \item Examine how gradients vary layer by layer.
  \item Discuss effective anisotropy in cyclically layered strata.
\end{itemize}

\paragraph{Connection}
\begin{itemize}
  \item Tie back to Week 2 physical-properties mixing demos.
  \item Discuss refraction in tilted and folded strata.
\end{itemize}

\timeline[slot]{Break}{10}

\timeline[item]{Energy Conservation and Geotherm Curvature}{35}
\paragraph{Central Question}
How can two terranes with identical surface heat flow have different Moho temperatures?

\paragraph{Steps}
\begin{itemize}
  \item Write the energy balance:
  \[
    q_0 = q_b + \int_0^{z_c} A(z)\,dz
  \]
  \item Infer relative crustal heat production given different basal inputs.
  \item Introduce the steady-state solution and the curvature term:
  \[
    T(z) = T_0 + \frac{q_0}{k}z - \frac{A}{2k}z^2
  \]
\end{itemize}

Curvature differences imply different deep temperatures even when surface gradients match.

\timeline[item]{Why Is Archean Heat Production So Low?}{35}
\paragraph{Observation}
Decay alone predicts high early heat production, yet preserved Archean crust shows limits.

\paragraph{Guided Discussion Topics}
\begin{itemize}
  \item Differentiation and partitioning of LILEs (K, Th, U).
  \item Selective erosion of felsic upper crust.
  \item Reworking: remelting and redistribution of older crust.
  \item Foundering of dense mafic roots.
  \item Secular mantle thermal evolution.
\end{itemize}

\timeline[slot]{Wrap-up}{20}
\begin{itemize}
  \item Summarize roles of $k$ and $A$.
  \item Connect observations to upcoming steady-state problem set.
\end{itemize}

\timeline[slot]{Preview: Transient Heat Flow}{10}

So far we have assumed that temperature does not change with time---a useful idealization for mature, stable terranes. But many geologically interesting situations involve heat flow that evolves: cooling of the oceanic lithosphere, heating around a magmatic intrusion, or the thermal relaxation of a thrust belt. Next week we move to transient heat conduction, governed by the diffusion equation, and develop the concept of thermal diffusivity as the quantity that controls how quickly temperature anomalies propagate through rock.

\workshopend
\notepage
\notepage

\subimport{activities/}{activity_refraction}
\subimport{activities/}{activity_layered_conductivity}
\subimport{activities/}{activity_conservation_of_energy}
\subimport{activities/}{activity_archean_heat_production}